%!TEX program = xelatex
%!TEX root = geometria_analitica.tex
%%Usar makeindex -s indexstyle.ist arquivo.idx no terminal para gerar o índice remissivo agrupado por inicial
%%Aós executar pdflatex arquivo

\chapter{Reta e Plano no Espaço}

\section{Equações do Plano}

Da Geometria Plana, sabemos que todo plano $\pi$ contém pelo menos três pontos não colineares, digamos $A$, $B$ e $C$.
\begin{figure}[H]
    \centering
    \includegraphics[scale=0.2]{pontos-plano.png}
\end{figure}


Assim considere os vetores
\begin{align*}
    \vec{u} = \vec{AB} \\
    \vec{v} = \vec{AC}.
\end{align*}
Como $\vec{u}$ e $\vec{v}$ não são paralelos, então podemos considerar o vetor $\vec{u}\times\vec{v}\ne \vec{0}$.
\begin{figure}[H]
    \centering
    \caption{Equação do plano}\label{definicao-plano}
    \includegraphics[scale=0.2]{equacao-plano.png}
\end{figure}
Assim, na Figura \ref{definicao-plano} vemos que um ponto $D$ pertence ao plano $\pi$ se, e somente se, $\vec{u}\times\vec{v} \perp \vec{AD}$. Logo um plano $\pi$ é um conjunto de vetores perpendiculares a um dado vetor.

Fixe então um ponto $P_0$ em $\real^3$ e $\vec{v}\ne \vec{0}$ um vetor. Passando por $P_0$ existe um \'unico plano $\pi$ perpendicular ao vetor $\vec{v}$. Assim um ponto $P$ do espaço pertence ao plano $\pi$ se, e somente se, $\vec{P_0P} \perp \vec{v}$, isto é,
\begin{equation}\label{equacaoplanovetorial}
    \inner{P_0P}{v} = 0.
\end{equation}
Sejam $\vec{v} = (a,b,c)$ um vetor, $P_0(x_0, y_0, z_0)$ e $P(x, y, z)$ pontos em $\real^3$. Logo $\vec{P_0P} = (x - x_0, y - y_0, z - z_0)$ e então da equação \eqref{equacaoplanovetorial} obtemos
\begin{equation}\label{equacaoplano}
    (x - x_0)a + (y - y_0)b + (z - z_0)c = 0.
\end{equation}

A equação \eqref{equacaoplano} é chamada de \textbf{equação cartesiana do plano} $\pi$. O vetor $\vec{v}$ é chamado de \textbf{vetor normal} ao plano $\pi$.\index{Plano!Equação Paramétrica}\index{Plano!Vetor Normal}

\begin{exemplos}
    Encontre a equação do plano $\pi$ nas seguintes situações:
    \begin{enumerate}
        \item O vetor normal é $\vec{v} = (1,2,-1)$ e o ponto de $\pi$ é $P_0(1,3,-1)$.
              \begin{solucao}
                  Se $P(x,y,z)$ pertence a $\pi$, então
                  \begin{align*}
                       & \inner{P_0P}{v} = 0                               \\
                       & \langle(x - 1, y - 3, z + 1), (1,2,-1)\rangle = 0 \\
                       & x - 1 + (y - 3)2 + (z + 1)(-1) = 0.
                  \end{align*}
                  Logo a equação do plano procurado é
                  \[
                      x + 2y - z - 8 = 0.
                  \]
              \end{solucao}
        \item O vetor normal é $\vec{v} = (0,0,1)$ e $P_0(0,0,0)$.
              \begin{solucao}
                  Se $P(x,y,z)$ pertence a $\pi$, então
                  \begin{align*}
                       & \inner{P_0P}{v} = 0                   \\
                       & \langle(x, y, z), (0,0,1)\rangle = 0.
                  \end{align*}
                  Logo a equação do plano procurado é
                  \[
                      z = 0.
                  \]
                  \begin{figure}[H]
                      \centering
                      \caption{Plano $z = 0$}
                      \includegraphics[scale=0.2]{plano-z=0.png}
                  \end{figure}
                  Agora, se $\vec{v} = (0,1,0)$ e $P_0(0,0,0)$, então a equação do plano será $y = 0$.
                  \begin{figure}[H]
                      \centering
                      \caption{Plano $y = 0$}
                      \includegraphics[scale=0.2]{plano-y=0.png}
                  \end{figure}
                  Agora, se $\vec{v} = (1,0,0)$ e $P_0(0,0,0)$, então a equação do plano será $x = 0$.
                  \begin{figure}[H]
                      \centering
                      \caption{Plano $x = 0$}
                      \includegraphics[scale=0.2]{plano-x=0.png}
                  \end{figure}
              \end{solucao}
        \item Passando pelos pontos $A(3,1,-2)$, $B(5,2,1)$ e $C(2,0,2)$.
              \begin{solucao}
                  Sejam
                  \begin{align*}
                      \vec{u} = \vec{AB} = (2,1,3) \\
                      \vec{v} = \vec{AC} = (-1,-1,4).
                  \end{align*}
                  Um vetor normal ao plano é dado por $\vec{w} = \vec{u}\times\vec{v}$. Assim
                  \[
                      \vec{w} = \det \begin{bmatrix}
                          \vec{i} & \vec{j} & \vec{k} \\
                          2       & 1       & 3       \\
                          -1      & -1      & 4
                      \end{bmatrix} = 4\vec{i} - 3\vec{j} - 2\vec{k} + \vec{k} + 3\vec{i} - 8\vec{j}.
                  \]
                  Assim $\vec{w} = (7,-11,-1)$ e portanto a equação do plano é
                  \[
                      7x - 11y - z - 12 = 0.
                  \]
              \end{solucao}
    \end{enumerate}
\end{exemplos}

Dada uma equação da forma
\[
    ax + by + cz + d = 0,
\]
onde $a$, $b$ e $c$ não são simultaneamente nulos, existe um plano $\pi$ representado por esta equação? A resposta é afirmativa. Suponha que $a \ne 0$. Tome $\vec{v} = (a,b,c) \ne 0$ e $P_0(-d/a,0,0)$. A equação do plano com vetor normal $\vec{v}$ e passando por $P_0$ é dada por
\begin{align*}
     & \inner{P_0P}{v} = 0                                          \\
     & \langle\left(x + \dfrac{d}{a},y,z\right), (a,b,c)\rangle = 0 \\
     & ax + by + cz + d = 0.
\end{align*}
Portanto qualquer equação da forma $ax + by + cz + d = 0$, com $a$, $b$ e $c$ não simultaneamente nulos define um plano.

Sejam $\pi_1$ e $\pi_2$ dois planos dados por
\begin{align*}
    \pi_1 : a_1x + b_1y + c_1z + d_1 = 0 \\
    \pi_2 : a_2x + b_2y + c_2z + d_2 = 0
\end{align*}
onde os vetores normais de $\pi_1$ e $\pi_2$ são, respectivamente,
\begin{align*}
    \vec{v_1} = (a_1, b_1, c_1) \\
    \vec{v_2} = (a_2, b_2, c_2).
\end{align*}
Os planos $\pi_1$ e $\pi_2$ são \textbf{paralelos} se, e somente se, $\vec{v_1} \varparallel\vec{v_2}$. Para que $\pi_1$ e $\pi_2$ sejam iguais devemos ter $\vec{v_1}\varparallel\vec{v_2}$  e dados $P_1\in\pi_1$ e $P_2\in\pi_2$ o vetor $\vec{P_1P_2}$ deve ser tal que $\vec{P_1P_2}\perp\vec{v_1}$ ou $\vec{P_1P_2}\perp\vec{v_2}$.

\begin{exemplo}
    Dados os seguintes planos, determinar quais são iguais ou paralelos dois a dois:
    \begin{center}
        \begin{tabular}{cc}
            $\pi_1 : 2x - y + 3z = 0$ &
            $\pi_2 : x - y - 3 = 0$                           \\
            $\pi_3 : y - 2x - 3z = 2$ & $\pi_4 : 2x - 2y = 6$ \\
            $\pi_5 : x = 0$.
        \end{tabular}
    \end{center}
    \begin{solucao}
        Os vetores normais a cada plano são
        \begin{center}
            \begin{tabular}{cc}
                $\vec{v_1} = (2, -1, 3)$  & $\vec{v_2} = (1,-1,0)$ \\
                $\vec{v_3} = (-2, 1, -3)$ & $\vec{v_2} = (2,-2,0)$ \\
                $\vec{v_5} = (1, 0, 0)$.
            \end{tabular}
        \end{center}
        Como $\vec{v_1} = -\vec{v_3}$ então $\vec{v_1}\varparallel\vec{v_2}$. Agora $P_1(0,0,0)\in\pi_1$ e $P_2(0,2,0)\in\pi_3$. Daí
        \[
            \inner{P_1P_2}{v_1} = \langle(0,2,0), (2,-1,3)\rangle = -2 \ne 0.
        \]
        Logo $\pi_1$ e $\pi_3$ são paralelos e distintos.
        Agora, como $\vec{v_4} = 2\vec{v_2}$ então $\vec{v_2}\varparallel\vec{v_4}$. Mas $P_1(3,0,0)\in\pi_2$ e $P_2(0,-3,0)\in\pi_4$. Daí
        \[
            \inner{P_1P_2}{v_2} = \langle(-3,-3,0), (1,-1,0)\rangle = 0.
        \]
        Logo $\pi_2 =\pi_4$.

        Os demais planos não são paralelos.
    \end{solucao}
\end{exemplo}

\section{Equações da Reta}

Seja $\vec{v}$ um vetor e $A$ um ponto. Sabemos que a reta $r$ passando por $A$ e paralela ao vetor $\vec{v}$ tem equação
\begin{equation}
    \vec{AP} = t\vec{v}
\end{equation}
onde $t \in \real$ e $P$ é um ponto qualquer de $r$. Sendo $A$ um ponto em $\real^3$ e $\vec{v}$ um vetor de $\real^3$ temos
\begin{align*}
    A = (x_0, y_0, z_0) \\
    \vec{v} = (a, b, c).
\end{align*}
Daí um ponto $P(x, y, z)$ pertence à reta $r$ se, e somente se,
\[
    \vec{AP} = (x - x_0, y - y_0, z - z_0) = t(a, b, c).
\]
Logo $P(x, y, z)$ pertence à reta $r$, se e só se, $P(x,y,z)$ satisfaz as equações
\begin{equation}\label{equacoesparametricasretaespaco}
    \begin{cases}
        x = x_0 + ta \\
        y = y_0 + tb \\
        z = z_0 + tc
    \end{cases}; \quad t \in \real.
\end{equation}
Tais equações são chamadas de \textbf{equações paramétricas} da reta $r$. O vetor $\vec{v}$ é chamado de \textbf{vetor diretor} de $r$.\index{Reta!Equações Paramétricas}\index{Reta!Vetor Diretor}

\begin{exemplo}
    As equações paramétricas da reta $r$ passando por $P_0(-3,3/2,4)$ e paralela ao vetor $\vec{v} = (-6,1,4)$ são
    \[
        \begin{cases}
            x = -3 - 6t \\
            y = 3/2 + t \\
            z = 4 + 4t
        \end{cases},\quad t \in \real.
    \]
\end{exemplo}

De forma análoga ao casa do plano, se $\vec{v} = (a,b,c)$ é tal que $abc \ne 0$, então podemos reescrever \eqref{equacoesparametricasretaespaco} como
\[
    \dfrac{x - x_0}{a} = \dfrac{y - y_0}{b} = \dfrac{z - z_0}{c}
\]
que são chamadas de \textbf{equações na forma simétrica} da reta $r$.\index{Reta!Equações Simétricas}

\section{Interseções}

\subsection{Entre retas}
Considere duas retas $r$ e $s$ de equações
\[
    r: \begin{cases}
        x = x_0 + ta \\
        y = y_0 + tb \\
        z = z_0 + tc
    \end{cases}; \quad s: \begin{cases}
        x = x_1 + \alpha d \\
        y = y_1 + \alpha e \\
        z = z_1 + \alpha f
    \end{cases}; t,\ \alpha \in \real.
\]
Os vetores diretores de $r$ e $s$ são $\vec{u} = (a,b,c)$ e $\vec{v} = (d,e,f)$ respectivamente.
\begin{enumerate}
    \item Se $\vec{u}\varparallel\vec{v}$, então $r$ e $s$ são paralelas ou coincidentes. Elas serão coincidentes se, e somente se, possuírem algum ponto em comum.\index{Retas!Paralelas}
          \begin{figure}[H]
              \centering
              \caption{Retas Paralelas}
              \includegraphics{retas-paralelas-espaco.pdf}
          \end{figure}
    \item Se os vetores diretores não são paralelos, temos duas possibilidades. Considere os pontos $P_0(x_0,y_0,z_0) \in r$ e $P_1(x_1,y_1,z_1) \in s$.
          \begin{enumerate}
              \item Se $\vec{P_1P_2}$, $\vec{u}$ e $\vec{v}$ estão no mesmo plano, isto é, $\langle\vec{P_1P_2}, \vec{u}\times\vec{v}\rangle = 0$ então $r$ e $s$ são concorrentes.\index{Retas!Concorrentes}
                    \begin{figure}[H]
                        \centering
                        \caption{Retas Concorrentes}
                        \includegraphics{retas-concorrentes-espaco.pdf}
                    \end{figure}
              \item Se $\vec{P_1P_2}$, $\vec{u}$ e $\vec{v}$ estão planos diferentes, isto é, $\langle\vec{P_1P_2}, \vec{u}\times\vec{v}\rangle \ne 0$ então $r$ e $s$ são chamadas de  \textbf{retas reversas}.\index{Retas! Reversas}
                    \begin{figure}[H]
                        \centering
                        \caption{Retas Reversas}
                        \includegraphics{retas-reversas-espaco.pdf}
                    \end{figure}
          \end{enumerate}
\end{enumerate}

\begin{exemplos}
    Determine se os seguintes pares de retas são paralelas, concorrentes ou reversas:
    \begin{enumerate}
        \item $r: \dfrac{x - 2}{2} = y + 3 = \dfrac{z - 2}{3}$; $s:\dfrac{x}{4} = \dfrac{y+3}{2} = \dfrac{z-2}{6}$
              \begin{solucao}
                  Como $\vec{u} = (2,1,3)$ e $\vec{v} = (4,2,6)$ são vetores diretores de $r$ e $s$ respectivamente, então $\vec{u}\varparallel\vec{v}$. Mas $P_0(0,-3,2)\in s$ e $P_0\notin r$, logo $r$ e $s$ são paralelas e distintas.
              \end{solucao}
        \item
              \[
                  r: \begin{cases}
                      x = 1 + 2t \\
                      y = 1 + 2t \\
                      z = 1 + t
                  \end{cases},\quad s: \begin{cases}
                      x = \alpha \\
                      y = \alpha \\
                      z = 0
                  \end{cases}
              \]
              \begin{solucao}
                  Como $\vec{u} = (2,2,1)$ e $\vec{v} = (1,1,0)$ são vetores diretores de $r$ e $s$ respectivamente, então $r$ e $s$ são concorrentes ou reversas. Igualando as equações de $r$ e $s$ obtemos
                  \begin{align*}
                      1 + 2t = \alpha \\
                      1 + 2t = \alpha \\
                      1 + t = 0
                  \end{align*}
                  donde $t = -1$ e $\alpha = -1$. Portanto $r$ e $s$ são concorrentes e o ponto de interseção é $P(-1,-1,0)$.
              \end{solucao}
        \item
              \[
                  r: \begin{cases}
                      x = 1 + 2t \\
                      y = 2 + t  \\
                      z = -1 + 2t
                  \end{cases},\quad s: \begin{cases}
                      x = 2 + \alpha \\
                      y = 1 + \alpha \\
                      z = -1 + 2\alpha
                  \end{cases}
              \]
              \begin{solucao}
                  Como $\vec{u} = (2,1,2)$ e $\vec{v} = (1,1,2)$ são vetores diretores de $r$ e $s$ respectivamente, então $r$ e $s$ são concorrentes ou reversas. Igualando as equações de $r$ e $s$ obtemos
                  \begin{equation}\label{exemploretasreversas}
                      \begin{cases}
                          1 + 2t = 2 + \alpha \\
                          2 + t = 1 + \alpha  \\
                          -1 + 2t = -1 + 2\alpha
                      \end{cases}
                  \end{equation}
                  Da primeira e segunda equação obtemos o sistema
                  \[
                      \begin{cases}
                          2t - \alpha = 1 \\
                          t - \alpha = -1
                      \end{cases}
                  \]
                  Cuja solução é $t = 2$ e $\alpha = 3$. Substituindo esse valores na terceira equação de \eqref{exemploretasreversas} obtemos $3 \ne 5$. Portanto $r$ e $s$ são reversas.
              \end{solucao}
    \end{enumerate}
\end{exemplos}

\subsection{Entre planos}

Dados três planos $\pi_1$, $\pi_2$ e $\pi_3$ de equações
\begin{align*}
    \pi_1 : a_1x + b_1y + c_1z + d_1 = 0 \\
    \pi_2 : a_2x + b_2y + c_2z + d_2 = 0 \\
    \pi_3 : a_3x + b_3y + c_3z + d_3 = 0.
\end{align*}
A interseção entre eles será constituída pelos pontos cujas coordenadas sejam as soluções do sistema formado pelas equações de $\pi_1$, $\pi_2$ e $\pi_3$ simultaneamente, ou seja, devem ser soluções do sistema
\[
    \begin{cases}
        a_1x + b_1y + c_1z + d_1 = 0 \\
        a_2x + b_2y + c_2z + d_2 = 0 \\
        a_3x + b_3y + c_3z + d_3 = 0.
    \end{cases}
\]

\begin{figure}[H]
    \centering
    \caption{$\pi_1\cap\pi_2\cap\pi_3 = \{P\}$}
    \includegraphics[scale=0.25]{intersecao-tres-planos.png}
\end{figure}

\begin{figure}[H]
    \caption{$\pi_1\cap\pi_2\cap\pi_3 = \emptyset$}
    \centering
    \includegraphics[scale=0.2]{tres_planos_dois_a_dois_retas.png}
\end{figure}

\begin{figure}[H]
    \centering
    \caption{$\pi_1\cap\pi_2\cap\pi_3 = \emptyset$}
    {\centering\includegraphics[scale=0.2]{tres-planos-paralelos.png}}
\end{figure}
\begin{exemplos}
    Encontre a interseção entre os planos:
    \begin{enumerate}
        \item $\pi_1 : x = y$, $\pi_2 : 2x - z = 0$ e $\pi_3 : 3x + 2y + z + 3 = 0$
              \begin{solucao}
                  A interseção será dada pela solução do sistema
                  \[
                      \begin{cases}
                          x - y = 0  \\
                          2x - z = 0 \\
                          3x + 2y + z + 3 = 0.
                      \end{cases}
                  \]
                  Resolvendo tal sistema obtemos que $\pi_1\cap\pi_2\cap\pi_3 = \{(-3/7,-3/7,-6/7)\}$.
              \end{solucao}
        \item $\pi_1 : x - 2y + 3z = 2$, $\pi_2 : 3x + y - 2z = 1$ e $\pi_3 : 5x - 3y + 4z = 4$
              \begin{solucao}
                  A interseção será dada pela solução do sistema
                  \[
                      \begin{cases}
                          x - 2y + 3z = 2 \\
                          3x + y - 2z = 1 \\
                          5x - 3y + 4z = 4
                      \end{cases}.
                  \]
                  Resolvendo tal sistema obtemos que $\pi_1 \cap \pi_2 \cap \pi_3 = \emptyset$. No entanto $\pi_i \cap \pi_j$ é uma reta para $i \ne j$.
              \end{solucao}
        \item $\pi_1 : 2x + y = 1$, $\pi_2 : x - y + z = 0$
              \begin{solucao}
                  A interseção será dada pela solução do sistema
                  \[
                      \begin{cases}
                          2x + y = 1 \\
                          x - y + z = 0.
                      \end{cases}
                  \]
                  Resolvendo tal sistema obtemos $z = 1 - 3x$ e $y = 1 - 2x$. Logo
                  \[
                      \pi_1\cap\pi_2 = \{(x,1 - 2x,1 - 3x) \mid x \in \real\}.
                  \]
                  Fazendo $x = t$ podemos escrever
                  \[
                      \begin{cases}
                          x = t      \\
                          y = 1 - 2t \\
                          z = 1 - 3t
                      \end{cases}; t \in \real
                  \]
                  que é a equação de uma reta passando pelo ponto $P_0(0,1,1)$ e com vetor diretor $\vec{v} = (1,-2,-3)$.
              \end{solucao}
    \end{enumerate}
\end{exemplos}

\subsection{Entre retas e planos}

Para encontrar a interseção de uma reta $r$ com um plano $\pi$, basta resolver o sistema formado pelas equações de $r$ e de $\pi$. Se o sistema tiver uma única solução $P_0(x_0, y_0, z_0)$, então $r \cap \pi = \{P_0\}$. Se o sistema for indeterminado então $r\subset \pi$ e se o sistema for impossível então $r\cap\pi = \emptyset$ e $r \varparallel \pi$.
\begin{figure}[H]
    \centering
    \caption{$r\cap\pi = \{P_0\}$}
    \includegraphics[scale=0.2]{intersecao-reta-plano.png}
\end{figure}

\begin{figure}[H]
    \centering
    \caption{$r\cap\pi= \emptyset$}
    \includegraphics[scale=0.3]{intersecao-reta-plano-vazia.png}
\end{figure}

\begin{figure}[H]
    \centering
    \caption{$r \subset \pi$}
    \includegraphics[scale=0.2]{reta-contida-no-plano.png}
\end{figure}

\begin{exemplos}
    Encontre a interseção entre a reta $r$ e o plano $\pi$ nos seguintes casos:
    \begin{enumerate}
        \item $r: x = t + 1$, $y = 1$, $z = 1$ e $\pi : 2x - y = 3$.
              \begin{solucao}
                  Substituindo as equações da reta na do plano obtemos $t = 1$. Logo $r\cap\pi = \{(2,1,1)\}$.
              \end{solucao}
        \item
              \[
                  r: \begin{cases}
                      2x - y + z = 0 \\
                      x + y = 1
                  \end{cases};
              \]
              $\pi : x - 2y + z = 5$.
              \begin{solucao}
                  O sistema
                  \[
                      \begin{cases}
                          2x - y + z = 0 \\
                          x + y = 1      \\
                          x - 2y + z = 5
                      \end{cases}
                  \]
                  não tem solução. Logo $r\cap\pi = \emptyset$.
              \end{solucao}
        \item
              \[
                  r: \begin{cases}
                      2x - y + z = 0 \\
                      x + y = 1
                  \end{cases};
              \]
              $\pi : x - y + 3z = -1$.
              \begin{solucao}
                  O sistema
                  \[
                      \begin{cases}
                          2x - y + z = 0 \\
                          x + y = 1      \\
                          x - y + 3z = -1
                      \end{cases}
                  \]
                  tem solução \'unica dada por $x = -3/7$ e $y = 10/7$ e $z = 2/7$. Logo $r\cap\pi = \{(-3/7, 10/7, 2/7)\}$.
              \end{solucao}
    \end{enumerate}
\end{exemplos}

\section{Distâncias}

\subsection{Entre ponto e plano}

Seja $P_1(x_1,y_1,z_1)$ um ponto e $\pi$ um plano de equação $ax + by + cz + d = 0$.
\begin{figure}[H]
    \centering
    \caption{Distância entre ponto e plano}
    \includegraphics[scale=0.2]{distancia-ponto-plano.png}
\end{figure}
Seja $r$ a reta que contém o ponto $P_1$ e é perpendicular ao plano $\pi$. Denote por $P_2(x_2,y_2,z_2)$ a interseção de $\pi$ com $r$. O ponto $P_2$ é chamado de \textbf{projeção ortogonal} de $P_1$ sobre $\pi$. A norma do vetor $\vec{P_1P_2}$ será a distância de $P_1$ a $\pi$, isto é,
\[
    d(P_1,\pi) = \norm{\vec{P_1P_2}}.
\]
Seja $\vec{u} = (a,b,c)$ o vetor normal de $\pi$. Como $\vec{u}\varparallel\vec{P_1P_2}$, existe $t \in \real$ tal que
\[
    \vec{P_1P_2} = t\vec{u}.
\]
Logo
\begin{equation}\label{equacaodistanciapontoplano}
    d(P_1,P_2) = \norm{t(a,b,c)} = |t|\sqrt{a^2 + b^2 + c^2}.
\end{equation}
Por outro lado, se $P(x_2, y_2, z_2)$, então
\[
    \vec{P_1P_2} = (x_2 - x_1, y_2 - y_1, z_2 - z_1) = t(a,b,c)
\]
e como $P_2 \in \pi$ devemos ter
\[
    a(x_1 + ta) + b(y_1 + tb) + c(z_1 + tc) + d = 0
\]
isto é,
\begin{equation}\label{equacaoauxiliardistanciapontoplano}
    t = -\dfrac{ax_1 + by_1 + cz_1 + d_1}{a^2 + b^2 + c^2}.
\end{equation}
Substituindo \eqref{equacaoauxiliardistanciapontoplano} em \eqref{equacaodistanciapontoplano} obtemos
\[
    d(P_1,\pi) = \left|\dfrac{ax_1 + by_1 + cz_1 + d_1}{a^2 + b^2 + c^2}\sqrt{a^2 + b^2 + c^2}\right|
\]
portanto
\[
    d(P_1,\pi) = \dfrac{|ax_1 + by_1 + cz_1 + d_1|}{\sqrt{a^2 + b^2 + c^2}}.
\]

Um outro modo de determinar essa distância é o seguinte: seja $P_0 \in \pi$. Então a distância de $P_1$ a $\pi$ é dada pela norma da projeção ortogonal do vetor $\vec{P_0P_1}$ na direção do vetor normal do plano $\vec{u} = (a,b,c)$. Assim
\[
    d(P_1,\pi) = \norm{\proj_{\vec{u}}\vec{P_0P_1}} = \dfrac{|\langle\vec{P_0P_1},\vec{u}\rangle|}{\norm{\vec{u}}}.
\]

\begin{figure}[H]
    \centering
    \caption{Distância entre ponto e plano}
    \includegraphics[scale=0.2]{distancia-ponto-plano-vetores.png}
\end{figure}

\newpage
\begin{exemplos}
    \begin{enumerate}
        \item A distância entre o plano $\pi: 2x + y - z = 4$ e o ponto $P(1,1,1)$ será
              \[
                  d(P,\pi) = \dfrac{|2 + 1 - 1 - 4|}{\sqrt{4 + 1 + 1}} = \dfrac{2}{\sqrt{6}}.
              \]
        \item Calcule a distância entre $\pi_1 : 2x - 3y + z - 2 = 0$ e $\pi_2 : 4x - 6y + 2z - 5 = 0$.
              \begin{solucao}
                  Um vetor normal de $\pi_1$ e $\vec{u} = (2,-3,1)$ e de $\pi_2$ é $\vec{v} = (4,-6,2)$. Como $\vec{u}\varparallel\vec{v}$, então $\pi_1\varparallel\pi_2$ e assim a distância entre $\pi_1$ e $\pi_2$ é dada pela distância de um ponto de $\pi_1$ a $\pi_2$, ou o contrário. Tomando $P_1(1,0,0)\in\pi_1$ temos
                  \[
                      d(\pi_1,\pi_2) = d(P_1,\pi_2) = \dfrac{|4-5|}{\sqrt{4^2 + (-6)^2 + 2^2}} = \dfrac{1}{\sqrt{56}}.
                  \]
              \end{solucao}
    \end{enumerate}
\end{exemplos}

\subsection{Entre um ponto e uma reta}

Seja $P$ um ponto e $r$ uma reta.
\begin{figure}[H]
    \centering
    \caption{Distância entre ponto e reta}
    \includegraphics[scale=0.3]{distancia-ponto-reta.png}
\end{figure}

Para determinar $d(P,r)$ primeiro construímos um plano $\pi$ passando por $P$ e com vetor normal paralelo ao vetor diretor de $r$. Assim
\[
    d(P,r) = d(P,P_1)
\]
onde $P_1 \in \pi\cap r$.

\begin{exemplo}
    Determinar a distância do ponto $P(1,2,-1)$ \`a reta
    \[
        r: \begin{cases}
            x = 1 + 2t \\
            y = 5 - t  \\
            z = -2 + 3t.
        \end{cases}
    \]
    \begin{solucao}
        Podemos tomar o vetor $\vec{u} = (2,1,-3)$ como vetor normal ao plano $\pi$ que contém $P$. Daí
        \begin{align*}
            \pi & : 2(x - 1) + (y - 2) - 3(z - 1) = 0 \\
            \pi & : 2x + y - 3z - 7 = 0.
        \end{align*}
        A interseção de $\pi$ com $r$ é dada por
        \[
            2(1 + 2t) + (5 - t) - 3(-2 + 3t) -7 = 0,
        \]
        isto é, $t = 1$ e então $P_1(3,4,1) \in \pi\cap r$. Portanto
        \[
            d(P,r) = d(P,P_1) = \sqrt{(1 - 3)^2 + (2 - 4)^2 + (-1 - 1)^2} = 2\sqrt{3}.
        \]
    \end{solucao}
\end{exemplo}

\subsection{Entre retas}

Sejam $r$ e $s$ retas. Se $r \varparallel s$ então a distância entre $r$ e $s$ é dada pela distância entre um ponto de $r$ e a reta $s$, ou o contrário. Suponha então que $r$ e $s$ são concorrentes ou reversas.

Estas duas retas definem dois planos paralelos (que podem ser coincidentes, no caso em que as retas são concorrentes). Um plano é o que contém $r$ e é paralelo a $s$, que vamos chamar de $\pi_1$. O outro contém $s$ e é paralelo a $r$, que chamaremos de $\pi_2$.

\begin{figure}[H]
    \centering
    \caption{Distância entre retas reversas}
    \includegraphics[scale=0.3]{distancia-retas-reversas-alternativa.png}
\end{figure}

Sejam $\vec{u}$ e $\vec{v}$ os vetores diretores de $r$ e $s$, respectivamente. Então o vetor $\vec{u}\times\vec{v}$ é normal aos planos $\pi_1$ e $\pi_2$. Assim
\[
    d(r,s) = d(\pi_1,\pi_2).
\]
Mas
\[
    d(\pi_1,\pi_2) = \dfrac{|\langle\vec{PQ},\vec{u}\times\vec{v}\rangle|}{\norm{\vec{u}\times\vec{v}}}
\]
onde $P \in r$ e $Q \in s$. Portanto,
\[
    d(r,s) = \dfrac{|\langle\vec{PQ},\vec{u}\times\vec{v}\rangle|}{\norm{\vec{u}\times\vec{v}}}.
\]

Se $r \cap s \ne \emptyset$, então $d(r,s) = 0$ pois $\vec{PQ}$, $\vec{u}$ e $\vec{v}$ estão no mesmo plano.

\begin{exemplo}
    Determine a distância entre as retas
    \[
        r: \begin{cases}
            x = 2 + t  \\
            y = 1 - 3t \\
            z = 1 + 2t
        \end{cases} \quad s: \begin{cases}
            x = -5 + 4\alpha \\
            y = 6 - 5\alpha  \\
            z = 4 + 3\alpha
        \end{cases}
    \]
    para $t$, $\alpha \in \real$.
    \begin{solucao}
        Os vetores diretores de $r$ e $s$ são $\vec{u} = (1,-3,2)$ e $\vec{v} = (4,-5,3)$, respectivamente. Logo $r$ e $s$ não são paralelas.
        % Agora, o sistema
        % \[
        %     \begin{cases}
        %         2 + t = -5 + 4s\\
        %         1 - 3t = 6 - 5s\\
        %         1 + 2t = 4 + 3s
        %     \end{cases}
        % \]
        % e assim $r$ e $s$ são reversas. Assim pelo ponto $(-5,6,4)$ de $s$ tracemos um reta $s'$ paralela a $r$. Daí
        % \[
        %     s' : \begin{cases}
        %         x = -5 + \alpha\\
        %         y = 6 - 3\alpha\\
        %         z = 4 + 2\alpha.
        %     \end{cases}
        % \]
        Temos $\vec{u}\times\vec{v} = (1,5,7)$. Tomando $P(2,1,1) \in r$ e $Q(-5,6,4) \in s$ temos $\vec{PQ} = (-7,5,3)$
        \[
            d(r,s) = = \dfrac{|\langle\vec{PQ},\vec{u}\times\vec{v}\rangle|}{\norm{\vec{u}\times\vec{v}}} = \dfrac{|-7+25+21|}{\sqrt{1 + 5^2 + 7^2}} = \dfrac{39}{\sqrt{75}}.
        \]

    \end{solucao}
\end{exemplo}
% subsection entre_retas (end)

% chapter reta_e_plano_no_espaco (end)
